% بخش: روش‌ها
% فصل: اثبات‌ها
% بخش فرعی: استفاده از تعریف‌ها

\documentclass[../../../include/open-logic-section]{subfiles}

\begin{document}

\olfileid{mth}{prf}{def}

\olsection{استفاده از تعریف‌ها}

گفتیم باید با همهٔ تعریف‌هایی که ممکن است در اثبات به کار روند آشنا
باشید و بتوانید آن‌ها را درست به کار ببرید. این نکته واقعاً مهم است
و ارزش دارد کمی دقیق‌تر بررسی شود. تعریف‌ها برای کوتاه‌کردن
ویژگی‌ها و روابط به کار می‌روند تا بتوانیم موجزتر درباره‌شان سخن
بگوییم. عبارت کوتاهِ معرفی‌شده را \emph{تعریف‌شونده} و آنچه این
عبارت کوتاه می‌کند \emph{تعریف‌کننده} می‌نامند. در اثبات‌ها اغلب
باید به چگونگی معرفیِ تعریف‌شونده بازگردیم، زیرا برای پیش‌برد اثبات
باید از ساختار منطقیِ تعریف‌کننده (صورت بلندی که اصطلاح تعریف‌شده کوتاه‌شدهٔ آن است)
بهره بگیریم. با بازکردن تعریف‌ها مطمئن می‌شوید که به کانون کنش
منطقی می‌رسید.

با یک مثال آغاز می‌کنیم. فرض کنید می‌خواهید حکم زیر را اثبات کنید:

\begin{prop}
برای هر دو مجموعهٔ $A$ و $B$، داریم $A \cup B = B \cup A$.
\end{prop}

حتی برای آغاز اثبات باید بدانیم یکسان‌بودن دو مجموعه یعنی چه؛ یعنی
باید بدانیم علامت «$=$» در این معادله برای مجموعه‌ها چه معنایی
دارد. مجموعه‌ها هنگامی یکسان تعریف می‌شوند که !!{element}های
یکسانی داشته باشند. پس تعریفی که باید باز کنیم چنین است:

\begin{defn}
مجموعه‌های $A$ و $B$ \emph{یکسان‌اند}، یعنی $A = B$، اگر و تنها
اگر هر !!{element} از~$A$ !!a{element} از~$B$ باشد و برعکس.
\end{defn}

این تعریف $A$ و~$B$ را به‌عنوان جانگهدارهایی برای مجموعه‌های دلخواه
به کار می‌برد. آنچه تعریف می‌کند — \emph{تعریف‌شونده} — عبارت
«$A = B$» است، و این کار را با دادن شرط صدقِ $A = B$ انجام می‌دهد.
این شرط — «هر !!{element} از~$A$ !!a{element} از~$B$ باشد و
برعکس» — \emph{تعریف‌کننده} است.\footnote{در این مورد خاص — و به‌شکلی
بسیار گیج‌کننده! — وقتی $A = B$، مجموعه‌های $A$ و $B$ فقط یک
مجموعهٔ واحدند، هرچند در دو سوی برابری حروف متفاوتی برای آن به کار
می‌بریم. اما شیوه‌های مشخص‌کردن آن مجموعه ممکن است متفاوت باشند، و
همین تعریف را نابدیهی می‌کند.} تعریف مشخص می‌کند که
$A = B$ درست است اگر، و تنها اگر (این را کوتاه می‌کنیم و
«اگر و تنها اگر» می‌گوییم)، شرط برقرار باشد.

هنگام به‌کاربردن تعریف باید $A$ و $B$ در تعریف را با موردی که
بررسی می‌کنید تطبیق دهید. در مورد ما، برای صدق
$A \cup B = B \cup A$ لازم است هر $z \in A \cup B$ در
$B \cup A$ نیز باشد و برعکس. عبارت $A \cup B$ در حکم نقشِ $A$ در
تعریف را بازی می‌کند و $B \cup A$ نقشِ~$B$ را. چون $A$ و $B$ هم
در تعریف و هم در حکم مورد اثبات به کار رفته‌اند، اما کاربردهایشان
متفاوت است، باید مراقب باشید آن دو کاربرد را با هم نیامیزید. برای
نمونه، اشتباه است گمان کنید با نشان‌دادن اینکه هر !!{element} از
$A$ !!a{element} از~$B$ است و برعکس، می‌توانید حکم را اثبات کنید؛
این کار $A = B$ را نشان می‌دهد، نه
$A \cup B = B \cup A$ را. (افزون بر این، چون $A$ و $B$ می‌توانند
هر دو مجموعه‌ای باشند، پیشرفت چندانی نخواهید کرد: وقتی چیزی دربارهٔ
$A$ و~$B$ فرض نشده باشد، چه‌بسا دو مجموعهٔ متفاوت باشند.)

در این اثبات با مفاهیم نظریه‌مجموعه‌ای مانند اجتماع سروکار داریم؛
پس برای فهم چگونگی ادامهٔ اثبات باید معنای نماد $\cup$ را نیز
بدانیم. گاهی بازکردن یک تعریف، تعریف‌های دیگری برای بازکردن پدید
می‌آورد. برای نمونه،
$A \cup B$ به‌صورت
$\Setabs{z}{z \in A \text{ یا } z \in B}$ تعریف می‌شود. پس اگر
می‌خواهید $x \in A \cup B$ را اثبات کنید، بازکردن تعریف $\cup$ به
شما می‌گوید باید
$x \in \Setabs{z}{z \in A \text{ یا } z \in B}$ را اثبات کنید.
اکنون باید به یاد آورید که
$x \in \Setabs{z}{\dots z\dots}$ اگر و تنها اگر $\dots x\dots$.
پس با بازکردن بیشترِ تعریف نماد
$\Setabs{z}{\dots z \dots}$، آنچه باید نشان دهید این است:
$x \in A$ یا $x \in B$. بنابراین، «هر !!{element} از
$A \cup B$ در $B \cup A$ نیز !!a{element} است» در واقع یعنی:
«برای هر $x$، اگر $x \in A$ یا $x \in B$، آن‌گاه
$x \in B$ یا $x \in A$». اگر تعریف‌های حکم را کاملاً باز کنیم،
می‌بینیم باید مطلب زیر را نشان دهیم:

\begin{prop}
برای هر دو مجموعهٔ $A$ و $B$: (الف) برای هر $x$، اگر
$x \in A$ یا $x \in B$، آن‌گاه $x \in B$ یا $x \in A$؛ و
(ب) برای هر $x$، اگر $x \in B$ یا $x \in A$، آن‌گاه
$x \in A$ یا $x \in B$.
\end{prop}

نکتهٔ مهم این است که بازکردن تعریف‌ها بخشی ضروری از ساختن اثبات
است. انجام درست آن گاهی دشوار است: باید مراقب باشید متغیرهای تعریف
و جمله‌های موجود در ادعای مورد اثبات را از هم بازشناسید و درست
تطبیق دهید. برای موفقیت باید بدانید پرسش چه می‌خواهد و همهٔ
اصطلاح‌های آن چه معنایی دارند؛ اغلب باید بیش از یک تعریف را باز
کنید. در اثبات‌های ساده‌ای مانند تمرین‌های زیر، راه‌حل تقریباً
بی‌درنگ از خود تعریف‌ها به دست می‌آید. البته همیشه به این سادگی
نخواهد بود.

\begin{prob}
فرض کنید از شما خواسته‌اند اثبات کنید
$A \cap B \neq \emptyset$. همهٔ تعریف‌های موجود در این عبارت را باز
کنید؛ یعنی آن را به صورتی بازگویید که در آن «$\cap$»، «=» یا
«$\emptyset$» ذکر نشود.
\end{prob}

\end{document}
